\section{结论}

本文提出了一种解决从2D图像恢复3D信息这一不适定逆问题的新范式。在该设定下，输入信号缺乏必要信息，模型必须借助从数据中学习到的先验才能做出有效预测。其他方法要么仅在2D空间中进行计算，要么使用额外的深度网络重建场景，而本文方法在3D空间中操作，并通过反向投影按需获取图像特征。本文方法有两个优势：（1）消除了中间层表示（如预测的深度图或点云），这些表示可能成为累积误差的来源；（2）通过将同一3D点投影到所有可用帧，利用多相机信息。

除了直接应用于自动驾驶3D目标检测外，本文工作还有几个值得未来探索的方向。例如，单点投影在检索到的图像特征图中产生有限的感受野，为每个目标查询采样多个点将为目标优化融入更多信息。此外，新检测头与输入模态无关，引入激光雷达（LiDAR）/雷达（RADAR）等其他模态将增强性能和鲁棒性。最后，将本文流程推广到室内导航和物体操作等其他领域，将扩大其应用范围并揭示更多改进方向。 